% 图 65.1 —— 由 `checks/emit_hand.py` 自动拼出（probe 精测 + prims2 图元分解）
%%
% ⚠ 这是**稿子不是成品**：跑 `checks/checkhand.py 65.1` 看指标 + 看叠合图。
%%
%% ⚠ 待核：
%   · 本图 probe 报出阴影族 1 个 —— **本脚本不生成阴影**（要照原书手写）；prims2 的 strokes 里可能含阴影线，核对时留意别画重
%   · 可疑字形 1 项（空心圈 ⊙ 常被认成字母 o）—— 见文件末
%%
\documentclass[12pt]{standalone}
\usepackage{fontspec}
\setsansfont{Arial}
\newfontfamily\mfallback{DejaVu Sans}
\usepackage{tikz}
\begin{document}
\begin{tikzpicture}[x=1pt, y=-1pt, line width=4.0pt, line cap=round, line join=round]

  %% ════ 直线段（prims2 的 strokes）════
  \draw[line width=4.0pt] (227.0,60.0) -- (227.0,17.0);
  \draw[line width=5.9pt] (298.0,143.0) -- (293.5,139.5);
  \draw[line width=5.0pt] (157.0,216.0) -- (217.0,216.0);
  \draw[line width=5.0pt] (757.0,217.0) -- (678.0,217.0);
  \draw[line width=5.0pt] (887.0,147.0) -- (871.9,147.0);
  \draw[line width=4.0pt] (156.0,217.0) -- (172.0,266.0);
  \draw[line width=5.0pt] (888.0,305.0) -- (887.0,376.0);
  \draw[line width=4.5pt] (888.0,210.0) -- (896.0,216.0);
  \draw[line width=4.0pt] (227.0,275.0) -- (227.0,225.0);
  \draw[line width=5.0pt] (934.0,150.0) -- (889.0,148.0);
  \draw[line width=4.0pt] (226.0,104.0) -- (227.0,206.0);
  \draw[line width=5.0pt] (1105.0,218.0) -- (1012.0,217.0);
  \draw[line width=5.0pt] (888.0,227.0) -- (887.0,301.0);
  \draw[line width=4.5pt] (322.0,215.0) -- (301.0,216.0);
  \draw[line width=6.0pt] (148.0,151.0) -- (150.0,142.0);
  \draw[line width=4.0pt] (148.0,151.0) -- (154.0,214.0);
  \draw[line width=5.0pt] (935.0,149.0) -- (934.0,138.3);
  \draw[line width=4.5pt] (999.0,125.0) -- (1004.0,136.0);
  \draw[line width=4.0pt] (953.0,217.0) -- (897.0,217.0);
  \draw[line width=5.0pt] (153.0,215.0) -- (129.0,215.0);
  \draw[line width=4.0pt] (889.0,60.0) -- (887.0,16.0);
  \draw[line width=5.0pt] (156.4,139.6) -- (151.0,141.0);
  \draw[line width=5.0pt] (1011.0,216.0) -- (1004.0,142.0);
  \draw[line width=5.0pt] (1010.0,217.0) -- (956.0,217.0);
  \draw[line width=4.0pt] (953.0,258.5) -- (955.0,241.3);
  \draw[line width=5.0pt] (955.0,241.3) -- (955.0,218.0);
  \draw[line width=5.0pt] (226.0,276.0) -- (181.0,273.0);
  \draw[line width=4.0pt] (890.0,61.0) -- (912.0,56.0);
  \draw[line width=5.0pt] (912.0,56.0) -- (934.8,53.0);
  \draw[line width=5.0pt] (934.8,53.0) -- (952.4,52.0);
  \draw[line width=5.0pt] (952.4,52.0) -- (969.6,54.0);
  \draw[line width=5.0pt] (888.0,146.0) -- (888.0,64.0);
  \draw[line width=4.0pt] (1004.0,137.0) -- (1004.0,141.0);
  \draw[line width=4.0pt] (227.0,384.0) -- (227.0,317.0);
  \draw[line width=5.0pt] (228.0,64.0) -- (227.0,100.0);
  \draw[line width=5.0pt] (954.0,216.0) -- (938.0,153.0);
  \draw[line width=5.0pt] (896.0,218.0) -- (889.0,226.0);
  \draw[line width=5.0pt] (226.0,224.0) -- (218.0,217.0);
  \draw[line width=3.0pt] (1072.0,69.0) -- (1085.0,69.0);
  \draw[line width=5.0pt] (325.0,216.0) -- (325.0,228.0);
  \draw[line width=5.0pt] (242.1,319.1) -- (228.0,317.0);
  \draw[line width=4.0pt] (806.0,217.0) -- (760.0,217.0);
  \draw[line width=5.0pt] (887.0,209.0) -- (888.0,149.0);
  \draw[line width=5.0pt] (125.0,215.0) -- (17.0,215.0);
  \draw[line width=4.5pt] (879.0,217.0) -- (886.0,210.0);
  \draw[line width=5.0pt] (879.0,217.0) -- (808.0,217.0);
  \draw[line width=5.0pt] (879.0,217.0) -- (887.0,226.0);
  \draw[line width=6.0pt] (172.0,267.0) -- (180.0,273.0);
  \draw[line width=5.0pt] (326.0,215.0) -- (439.8,216.0);
  \draw[line width=3.8pt] (439.8,216.0) -- (443.0,214.0);
  \draw[line width=5.0pt] (228.0,224.0) -- (235.0,216.0);
  \draw[line width=4.0pt] (886.0,387.0) -- (888.0,379.0);
  \draw[line width=6.0pt] (305.0,170.0) -- (299.0,144.0);
  \draw[line width=5.0pt] (228.0,207.0) -- (235.0,216.0);
  \draw[line width=5.0pt] (304.0,172.0) -- (300.0,215.0);
  \draw[line width=5.0pt] (886.0,377.0) -- (847.9,370.0);
  \draw[line width=5.0pt] (297.0,216.0) -- (235.0,216.0);
  \draw[line width=5.0pt] (227.0,277.0) -- (226.0,312.0);

  %% ════ 曲线（prims2 的 curves —— 三段式，坐标同为 (y,x)）════
  \draw[line width=5.0pt] (293.5,139.5) .. controls (284.1,112.3) and (257.9,95.2) .. (228.0,101.0);
  \draw[line width=4.0pt] (226.0,61.0) .. controls (179.8,53.3) and (146.2,97.4) .. (150.0,140.0);
  \draw[line width=5.0pt] (871.9,147.0) .. controls (832.4,141.6) and (802.6,178.4) .. (807.0,216.0);
  \draw[line width=5.0pt] (1011.0,218.0) .. controls (1012.8,289.5) and (976.1,389.1) .. (889.0,378.0);
  \draw[line width=4.0pt] (299.0,142.0) .. controls (295.4,105.9) and (265.6,70.3) .. (229.0,64.0);
  \draw[line width=4.0pt] (180.0,274.0) .. controls (188.3,292.0) and (206.0,307.6) .. (225.0,313.0);
  \draw[line width=5.0pt] (1003.0,142.0) .. controls (983.8,152.7) and (960.3,150.6) .. (939.0,152.0);
  \draw[line width=4.0pt] (148.0,151.0) .. controls (129.2,164.7) and (119.8,191.0) .. (126.0,214.0);
  \draw[line width=4.0pt] (934.0,138.3) .. controls (926.8,118.9) and (941.4,96.9) .. (960.9,93.1);
  \draw[line width=5.0pt] (960.9,93.1) .. controls (980.7,89.4) and (997.3,107.0) .. (999.0,125.0);
  \draw[line width=5.0pt] (225.0,103.0) .. controls (198.7,105.9) and (176.1,121.7) .. (156.4,139.6);
  \draw[line width=5.0pt] (889.0,304.0) .. controls (918.9,310.3) and (952.5,291.0) .. (953.0,258.5);
  \draw[line width=5.0pt] (969.6,54.0) .. controls (1006.7,56.3) and (1036.2,110.0) .. (1005.0,136.0);
  \draw[line width=5.0pt] (323.0,214.0) .. controls (319.0,199.6) and (316.9,181.0) .. (306.0,171.0);
  \draw[line width=4.0pt] (228.0,276.0) .. controls (261.5,276.7) and (294.2,250.8) .. (298.0,217.0);
  \draw[line width=4.0pt] (886.0,302.0) .. controls (845.2,296.4) and (807.9,260.6) .. (807.0,218.0);
  \draw[line width=4.0pt] (325.0,228.0) .. controls (336.6,276.4) and (295.8,338.3) .. (242.1,319.1);
  \draw[line width=4.0pt] (226.0,207.0) .. controls (222.0,206.9) and (217.9,211.0) .. (218.0,215.0);
  \draw[line width=5.0pt] (128.0,216.0) .. controls (130.4,238.7) and (150.0,260.0) .. (171.0,267.0);
  \draw[line width=5.0pt] (759.0,216.0) .. controls (765.8,146.6) and (818.3,81.3) .. (887.0,63.0);
  \draw[line width=5.0pt] (847.9,370.0) .. controls (781.6,354.7) and (754.4,280.7) .. (758.0,218.0);

  %% ════ 标签（probe 直接给的 \node 行）════
  %% ⚠ 全部补了 fill=white：文字在最上层，否则与曲线交叉干扰
     \node[anchor=center,inner sep=0pt,fill=white,font=\fontsize{30.7}{38.4}\selectfont] at (423.3,65.0) {\textsf{\textbf{2}}};   % box(415,54,15,21)
     \node[anchor=center,inner sep=0pt,fill=white,font=\fontsize{30.3}{37.9}\selectfont] at (1112.3,66.3) {\textsf{\textbf{9}}};   % box(1103,55,15,21)
     \node[anchor=center,inner sep=0pt,fill=white,font=\fontsize{52.0}{65.0}\selectfont] at (405.2,69.5) {\textsf{\textbf{|}}};   % box(396,56,15,20)
     \node[anchor=center,inner sep=0pt,fill=white,font=\fontsize{30.9}{38.7}\selectfont] at (342.8,69.3) {\textsf{\textbf{\textit{n}}}};   % box(334,59,16,17)
     \node[anchor=center,inner sep=0pt,fill=white,font=\fontsize{30.3}{37.9}\selectfont] at (379.8,67.6) {{\mfallback\textbf{−}}};   % box(364,61,17,4)
     \node[anchor=center,inner sep=0pt,fill=white,font=\fontsize{30.0}{37.5}\selectfont] at (1049.8,70.8) {\textsf{\textbf{\textit{n}}}};   % box(1041,61,16,16)
     \node[anchor=center,inner sep=0pt,fill=white,font=\fontsize{31.2}{39.0}\selectfont] at (1086.5,68.7) {{\mfallback\textbf{−}}};   % box(1070,62,18,4)
     \node[anchor=center,inner sep=0pt,fill=white,font=\fontsize{31.2}{39.0}\selectfont] at (379.5,73.7) {{\mfallback\textbf{−}}};   % box(363,67,18,4)

  % 钉死画布：否则超出的图元/控制点会把页面撑大，1:1 回比就失准
  \pgfresetboundingbox
  \path[use as bounding box] (0,0) rectangle (1133,398);
\end{tikzpicture}
\end{document}

% ⚠ 待核清单（probe 拿不准，人眼过一遍）：
%   · 未识别/跳过：⑤ 跳过/未识别的字形
